\section{An explicit hyperbolic punctured torus}
\label{app:punctured-torus}

For completeness, we recall the model in
\cite[Section~5.5]{Bonahon2009} that supplies the surface used in
Section~\ref{sec:cometric}.  In the upper half-plane, take the ideal
quadrilateral with vertices $-1,0,1,\infty$.  Pair its opposite sides by
\[
 \phi(z)=\frac{z+1}{z+2},\qquad
 \psi(z)=\frac{z-1}{2-z}.
\]
Both maps are orientation-preserving hyperbolic isometries:
$\phi$ sends the side with endpoints $-1,\infty$ to that with endpoints
$0,1$, and $\psi$ sends the side with endpoints $1,\infty$ to that with
endpoints $0,-1$.  Their side pairings identify all four ideal vertices.
The quotient has one face, two edges and one removed vertex; its
compactification is an oriented torus.

For $a>1$, choose the horodisk $\{\operatorname{Im}z\ge a\}$ at infinity
and horodisks of Euclidean radius $1/(2a)$ at $-1,0,1$.
These horodisks are disjoint and their boundary arcs are matched by the
side pairings.  The vertex cycle is represented, up to sign, by
\[
 \psi^{-1}\phi^{-1}\psi\phi
 =\begin{pmatrix}1&6\\0&1\end{pmatrix};
\]
thus the end unfolds to a horodisk modulo $z\mapsto z+6$.
Writing $\theta=x/6$ and $r=\log(y/a)$ gives
\[
 \frac{dx^2+dy^2}{y^2}
 =dr^2+\left(\frac6a\right)^2e^{-2r}d\theta^2,
 \qquad \theta\in\mathbb R/\mathbb Z.
\]
The side identifications are smooth along their geodesic interiors,
the displayed cusp is complete, and its complement in the quotient is
compact.  This gives the complete finite-area punctured torus and the
cusp normal form used in the main construction.
